% SEGMENT OLP-0058-S01
% Part: propositional-logic
% Chapter: syntax-and-semantics
% Section: formulas

\documentclass[../../../include/open-logic-section]{subfiles}

\begin{document}

\olfileid{pl}{syn}{fml}

\olsection{கூற்றுத் தருக்க \usetoken{P}{formula}}

கூற்றுத் தருக்கவியலின் !!^{formula}s,
\emph{!!{propositional
    variable}s}
\iftag{prvFalse}{\iftag{prvTrue}{,}{ மற்றும்} பொய்மைக்கான
  கூற்று மாறிலி~$\lfalse$}{}
\iftag{prvTrue}{\iftag{prvFalse}{
    மற்றும்}{} மெய்மைக்கான கூற்று மாறிலி~$\ltrue$}{} ஆகியவற்றிலிருந்து
\emph{தருக்க இணைப்புகளைப்} பயன்படுத்திக் கட்டமைக்கப்படுகின்றன.

\begin{enumerate}
\item !!^a{denumerable} கணமான~$\PVar$ இல் உள்ள !!{propositional variable}s
  $\Obj p_0$, $\Obj p_1$, \dots
\tagitem{prvFalse}{!!{falsity}க்கான கூற்று மாறிலி~$\lfalse$.}{}
\tagitem{prvTrue}{!!{truth}க்கான கூற்று மாறிலி~$\ltrue$.}{}
\item தருக்க இணைப்புகள்:
  \startycommalist
  \iftag{prvNot}{\ycomma $\lnot$ (மறுப்பு)}{}%
  \iftag{prvAnd}{\ycomma $\land$ (இணையல்)}{}%
  \iftag{prvOr}{\ycomma $\lor$ (பிரிப்பிணைவு)}{}%
  \iftag{prvIf}{\ycomma $\lif$ (!!{conditional})}{}%
  \iftag{prvIff}{\ycomma $\liff$ (!!{biconditional})}{}%
\item குறியிடல் குறிகள்: (, ), மற்றும் காற்புள்ளி.
\end{enumerate}

கூற்றுத் தருக்கவியலின் இந்த மொழியை $\Lang L_0$ என்று குறிக்கிறோம்.

% SEGMENT OLP-0058-S02
\iftag{defNot,defOr,defAnd,defIf,defIff,defTrue,defFalse,defEx,defAll}{%
மேலே அறிமுகப்படுத்தப்பட்ட மூல இணைப்புகளுடன், பின்வரும்
\emph{வரையறுக்கப்பட்ட} குறியீடுகளையும் பயன்படுத்துகிறோம்:
\startycommalist
  \iftag{defNot}{\ycomma $\lnot$ (மறுப்பு)}{}%
  \iftag{defAnd}{\ycomma $\land$ (இணையல்)}{}%
  \iftag{defOr}{\ycomma $\lor$ (பிரிப்பிணைவு)}{}%
  \iftag{defIf}{\ycomma $\lif$ (!!{conditional})}{}%
  \iftag{defIff}{\ycomma $\liff$ (!!{biconditional})}{}%
  \iftag{defFalse}{\ycomma $\lfalse$ (!!{falsity})}{}%
  \iftag{defTrue}{\ycomma $\ltrue$ (!!{truth})}}{}.

\begin{tagblock}{defNot,defOr,defAnd,defIf,defIff,defTrue,defFalse,defEx,defAll}
\begin{explain}
வரையறுக்கப்பட்ட குறியீடு அதிகாரப்பூர்வமாக மொழியின் பகுதியன்று; அது ஒரு
முறையற்ற சுருக்கக்குறியாக அறிமுகப்படுத்தப்படுகிறது. மூலக் குறியீடுகளை
மட்டுமே பயன்படுத்தினால் மிகவும் நீளமாகிவிடும் வாய்பாடுகளைச் சுருக்கமாக
எழுத அது உதவுகிறது. இது தெளிவான நன்மை. ஆனால் நிறுவல்கள் சுருங்குவது
இதைவிடப் பெரிய நன்மை. ஒரு குறியீடு மூலக் குறியீடாக இருந்தால், நிறுவல்களில்
அதைத் தனியாகக் கையாள வேண்டும். எனவே மூலக் குறியீடுகள் அதிகமானால் நமது
நிறுவல்களும் நீளமாகும்.
\end{explain}
\end{tagblock}

% SEGMENT OLP-0058-S03
% Alternate symbols

\begin{intro}
மேலே நாம் பயன்படுத்தியவற்றிலிருந்து வேறுபட்ட கலைச்சொற்களையும்
குறியீடுகளையும் நீங்கள் அறிந்திருக்கலாம். தருக்கவியல் நூல்களும்
ஆசிரியர்களும் ``மறுப்பு'' என்பதற்கு பொதுவாக $\sim$, $\neg$, மற்றும்~!\
ஆகியவற்றையும், ``இணையல்'' என்பதற்கு $\wedge$, $\cdot$, மற்றும் $\&$
ஆகியவற்றையும் பயன்படுத்துகின்றனர். ``நிபந்தனைக்கூற்று'' அல்லது
``உட்குறிப்பு'' என்பதற்குப் பொதுவாகப் பயன்படுத்தப்படும் குறியீடுகள்
$\rightarrow$, $\Rightarrow$, மற்றும் $\supset$ ஆகும்.
\iftag{prvIff,defIff}{``இரு நிபந்தனைக்கூற்று,'' ``இருவழி உட்குறிப்பு,''
  அல்லது ``(பொருட்சார்) சமானம்'' என்பதற்கான குறியீடுகள்
  $\leftrightarrow$, $\Leftrightarrow$, மற்றும் $\equiv$ ஆகும்.}{}
\iftag{prvFalse,defFalse}{$\lfalse$ குறியீடு ``பொய்மை,'' ``பொய்யம்,''
  ``முரண்பாடு,'' அல்லது ``கீழ்'' என்றும் பலவாறு அழைக்கப்படுகிறது.}{}
\iftag{prvTrue,defTrue}{$\ltrue$ குறியீடு ``மெய்மை,'' ``மெய்யம்,''
  அல்லது ``மேல்'' என்றும் பலவாறு அழைக்கப்படுகிறது.}{}
\end{intro}

% SEGMENT OLP-0058-S04
\begin{defn}[வாய்பாடு]
\ollabel{defn:formulas}
கூற்றுத் தருக்கவியலின் \emph{!!{formula}s} கணம்~$\Frm[L_0]$ பின்வருமாறு
தொகுத்தறிதலாக வரையறுக்கப்படுகிறது:
\begin{enumerate}
\tagitem{prvFalse}{$\lfalse$ ஓர் அணு !!{formula}.}{}

\tagitem{prvTrue}{$\ltrue$ ஓர் அணு !!{formula}.}{}

\item ஒவ்வொரு !!{propositional variable}~$\Obj p_i$ உம் ஓர் அணு
  !!{formula}.

\tagitem{prvNot}{$!A$ !!a{formula} எனில், $\lnot !A$ உம்
  !!a{formula}.}{}

\tagitem{prvAnd}{$!A$, $!B$ ஆகியவை !!{formula}s எனில், $(!A \land
  !B)$ உம் !!a{formula}.}{}

\tagitem{prvOr}{$!A$, $!B$ ஆகியவை !!{formula}s எனில், $(!A \lor !B)$
  உம் !!a{formula}.}{}

\tagitem{prvIf}{$!A$, $!B$ ஆகியவை !!{formula}s எனில், $(!A \lif !B)$
  உம் !!a{formula}.}{}

\tagitem{prvIff}{$!A$, $!B$ ஆகியவை !!{formula}s எனில், $(!A \liff !B)$
  உம் !!a{formula}.}{}

\tagitem{limitClause}{வேறெதுவும் !!a{formula} அன்று.}{}
\end{enumerate}
\end{defn}

% SEGMENT OLP-0058-S05
\begin{explain}
!!{formula}s இன் வரையறை ஒரு \emph{தொகுத்தறிதல் வரையறை}. அடிப்படையில்,
!!{formula}s இன் கணத்தை முடிவிலாப் பல நிலைகளில் கட்டமைக்கிறோம். தொடக்க
நிலையில் எல்லா அணு வாய்பாடுகளையும் வாய்பாடுகளாக அறிவிக்கிறோம்; இது
வரையறையின் முதல் சில நேர்வுகளான
\iftag{prvTrue}{$\ltrue$, }{}%
\iftag{prvFalse}{$\lfalse$, }{}%
$\Obj p_i$ ஆகிய நேர்வுகளுக்கு ஒத்தது. ஆகவே ``அணு !!{formula}'' என்பது
இந்த வடிவங்களில் ஏதேனும் ஒன்றில் உள்ள !!{formula}வைக் குறிக்கும்.

வரையறையின் பிற நேர்வுகள், ஏற்கெனவே கட்டமைக்கப்பட்ட !!{formula}s
இலிருந்து புதிய !!{formula}s ஐக் கட்டமைப்பதற்கான விதிகளைத் தருகின்றன.
இரண்டாம் நிலையில் அவற்றைப் பயன்படுத்தி அணு !!{formula}s இலிருந்து
!!{formula}s ஐக் கட்டமைக்கலாம். மூன்றாம் நிலையில் அணு வாய்பாடுகளிலிருந்தும்
இரண்டாம் நிலையில் பெறப்பட்டவற்றிலிருந்தும் புதிய வாய்பாடுகளைக்
கட்டமைக்கிறோம்; இவ்வாறு தொடர்கிறோம். இறுதியில் அத்தகைய ஏதோ ஒரு நிலையில்
கட்டமைக்கப்படுவது மட்டுமே !!{formula} ஆகும்.
\end{explain}

% SEGMENT OLP-0058-S06
$!B$, $!C$ ஆகியவற்றிலிருந்து இரு இட இணைப்பு~$\ast$ ஐப் பயன்படுத்திக்
கட்டமைக்கப்பட்ட வாய்பாடு $(!B \ast !C)$ ஐ எழுதும்போது, பெரும்பாலும்
வெளிப்புற அடைப்புக்குறிச் சோடியை விட்டுவிட்டு $!B \ast !C$ என்று மட்டும்
எழுதுவோம்.

% SEGMENT OLP-0058-S07
\begin{tagblock}{defNot,defOr,defAnd,defIf,defIff,defTrue,defFalse,defEx,defAll}
\begin{defn}
வரையறுக்கப்பட்ட செயலிகளைப் பயன்படுத்திக் கட்டமைக்கப்படும் வாய்பாடுகளைப்
பின்வருமாறு புரிந்துகொள்ள வேண்டும்:

\begin{tagenumerate}{defTrue,defFalse,defNot,defOr,defAnd,defIf,defIff,defEx,defAll}

\tagitem{defTrue}{$\ltrue$ என்பது
  \iftag{prvFalse}{$\lnot\lfalse$}{ஏதோ ஒரு நிலையான அணு
    !!{formula}~$!A$ க்கான $(!A \lor \lnot !A)$} என்பதன் சுருக்கம்.}{}

\tagitem{defFalse}{$\lfalse$ என்பது
  \iftag{prvTrue}{$\lnot\ltrue$}{ஏதோ ஒரு நிலையான அணு
    !!{formula}~$!A$ க்கான $(!A \land \lnot !A)$} என்பதன் சுருக்கம்.}{}

\tagitem{defNot}{$\lnot !A$ என்பது $!A \lif \lfalse$ என்பதன் சுருக்கம்.}{}

\tagitem{defOr}{$!A \lor !B$ என்பது
  \iftag{prvAnd}{$\lnot(\lnot !A \land \lnot !B)$}{$\lnot !A \lif
    !B$} என்பதன் சுருக்கம்.}{}

\tagitem{defAnd}{$!A \land !B$ என்பது
  \iftag{prvOr}{$\lnot(\lnot !A \lor \lnot !B)$}{$\lnot (!A \lif
    \lnot !B)$} என்பதன் சுருக்கம்.}{}

\tagitem{defIf}{$!A \lif !B$ என்பது
  \iftag{prvOr}{$\lnot !A \lor !B)$}{$\lnot (!A \land \lnot !B)$}
  என்பதன் சுருக்கம்.}{}

\tagitem{defIff}{$!A \liff !B$ என்பது $(!A \lif !B) \land (!B
  \lif !A)$ என்பதன் சுருக்கம்.}{}
\end{tagenumerate}
\end{defn}
\end{tagblock}

% SEGMENT OLP-0058-S08
\begin{defn}[தொடரமைப்புசார் ஒரேதன்மை]
$\ident$ குறியீடு குறியீட்டுச் சரங்களுக்கு இடையிலான தொடரமைப்புசார்
ஒரேதன்மையைக் குறிக்கிறது. அதாவது $!A \ident !B$ என்பது, அப்போதும்
அப்போது மட்டுமே, $!A$, $!B$ ஆகியவை ஒரே நீளமுள்ள குறியீட்டுச் சரங்களாகவும்
ஒவ்வொரு இடத்திலும் ஒரே குறியீட்டைக் கொண்டவையாகவும் இருப்பதைக் குறிக்கும்.
\end{defn}

தொடரிணைப்பால் பெறப்பட்ட சரங்கள் $\ident$ குறியீட்டின் இருபுறமும் வரலாம்.
எடுத்துக்காட்டாக, $!A \ident (!B \lor !C)$ என்பது: தொடக்க அடைப்புக்குறி,
$!B$ சரம், $\lor$ குறியீடு, $!C$ சரம், நிறைவு அடைப்புக்குறி ஆகியவற்றை
இந்த வரிசையில் தொடரிணைத்துப் பெறப்படும் சரமும் $!A$ குறியீட்டுச் சரமும்
ஒரே சரம் என்று பொருள்படும். இது உண்மையானால், $!A$ இன் முதல் குறியீடு
தொடக்க அடைப்புக்குறி என்றும், $!A$ ஆனது $!B$ ஐ ஓர் உட்சரமாகக்
கொண்டுள்ளது என்றும் (இரண்டாம் குறியீட்டிலிருந்து தொடங்கி), அந்த உட்சரத்தை
$\lor$ தொடர்ந்து வருகிறது என்றும், இன்ன பிறவும் அறிகிறோம்.

\end{document}
